\documentclass[11pt]{article}

\usepackage{amsmath,amssymb,amsthm,mathtools}
\usepackage[a4paper,margin=1in]{geometry}
\usepackage{enumitem}
\usepackage{bm}

\newtheorem{definition}{Definition}
\newtheorem{theorem}{Theorem}
\newtheorem{proposition}{Proposition}
\newtheorem{lemma}{Lemma}
\newtheorem{corollary}{Corollary}
\theoremstyle{remark}
\newtheorem{remark}{Remark}
\newtheorem{example}{Example}

\usepackage{tikz-cd}

\newcommand{\F}{\mathbb{F}}
\newcommand{\C}{\mathcal{C}}
\newcommand{\Cp}{\mathcal{C}^{\perp}}
\newcommand{\Cstar}{\mathcal{C}^{*}}
\newcommand{\ann}{^{\circ}}
\newcommand{\ip}[2]{\left\langle #1,#2\right\rangle}
\newcommand{\Hom}{\operatorname{Hom}}
\newcommand{\Ker}{\operatorname{Ker}}
\newcommand{\im}{\operatorname{Im}}
\newcommand{\Span}{\operatorname{span}}
\newcommand{\Null}{\operatorname{Null}}
\newcommand{\Row}{\operatorname{Row}}
\newcommand{\wt}{\operatorname{wt}}
\newcommand{\id}{\operatorname{id}}

\title{Lecture Notes on Dual Codes, Dual Spaces, Orthogonal Complements, and Quotient Interpretations}
\author{}
\date{}

\begin{document}
	
	\maketitle
	
	\tableofcontents
	
	\section{Introduction}
	
	These notes explain the relation among the following notions:
	
	\begin{itemize}[leftmargin=2em]
		\item the dual space \(V^* = \Hom(V,\F)\),
		\item the orthogonal complement \(W^\perp\) of a subspace \(W \subseteq V\),
		\item the dual code \(\C^\perp\) of a linear code \(\C \subseteq \F_q^n\),
		\item the quotient spaces \(V/W^\perp\), \(\F_q^n/\C^\perp\), and \(\F_q^n/\C\),
		\item the coding-theoretic meaning of generator matrices, parity-check matrices, and syndromes.
	\end{itemize}
	
	A central theme is that coding theory is linear algebra over a finite field. Many constructions in coding theory are best understood as special cases of standard constructions from bilinear algebra.
	
	\section{Dual Spaces in Linear Algebra}
	
	\begin{definition}
		Let \(V\) be a vector space over a field \(\F\). The \emph{dual space} of \(V\) is
		\[
		V^* := \Hom(V,\F),
		\]
		the vector space of all linear maps \(f : V \to \F\).
	\end{definition}
	
	Thus an element of \(V^*\) is not a vector in \(V\), but a linear functional on \(V\).
	
	\begin{example}
		If \(V=\F^k\), then every \(f\in V^*\) is of the form
		\[
		f(x_1,\dots,x_k)=a_1x_1+\cdots+a_kx_k
		\]
		for unique \(a_1,\dots,a_k\in\F\). Hence \(V^*\cong \F^k\), but this is an isomorphism of vector spaces, not equality as sets.
	\end{example}
	
	\begin{theorem}
		If \(V\) is finite-dimensional, then
		\[
		\dim V^*=\dim V.
		\]
	\end{theorem}
	
	\begin{proof}
		Choose a basis \(e_1,\dots,e_n\) of \(V\). Define \(e_1^*,\dots,e_n^*\in V^*\) by
		\[
		e_i^*(e_j)=\delta_{ij}.
		\]
		Then \(e_1^*,\dots,e_n^*\) form a basis of \(V^*\).
	\end{proof}
	
	\section{Bilinear Forms and Orthogonal Complements}
	
	\begin{definition}
		A \emph{bilinear form} on a vector space \(V\) is a map
		\[
		B:V\times V\to \F
		\]
		that is linear in each variable.
	\end{definition}
	
	Often we write
	\[
	\ip{x}{y}=B(x,y).
	\]
	
	\begin{definition}
		Let \(W\subseteq V\) be a subspace. The \emph{orthogonal complement} of \(W\) with respect to \(B\) is
		\[
		W^\perp := \{v\in V:\ip{w}{v}=0 \text{ for all } w\in W\}.
		\]
	\end{definition}
	
	\begin{proposition}
		\(W^\perp\) is a subspace of \(V\).
	\end{proposition}
	
	\begin{proof}
		If \(v_1,v_2\in W^\perp\) and \(a,b\in\F\), then for all \(w\in W\),
		\[
		\ip{w}{av_1+bv_2}=a\ip{w}{v_1}+b\ip{w}{v_2}=0.
		\]
		Hence \(av_1+bv_2\in W^\perp\).
	\end{proof}
	
	\begin{definition}
		A bilinear form \(B\) is \emph{nondegenerate} if
		\[
		\ip{x}{v}=0 \text{ for all } x\in V
		\quad\Longrightarrow\quad
		v=0.
		\]
	\end{definition}
	
	\begin{remark}
		For \(V=\F_q^n\), the standard bilinear form is
		\[
		\ip{x}{y}=x_1y_1+\cdots+x_ny_n.
		\]
		This is nondegenerate.
	\end{remark}
	
	\section{From a Bilinear Form to the Dual Space}
	
	Suppose \(V\) has a bilinear form \(B\). Then every vector \(v\in V\) defines a linear functional
	\[
	\varphi_v : V \to \F,\qquad \varphi_v(x)=\ip{x}{v}.
	\]
	This gives a linear map
	\[
	\Phi:V\to V^*,\qquad v\mapsto \varphi_v.
	\]
	
	\begin{proposition}
		If \(B\) is nondegenerate and \(V\) is finite-dimensional, then \(\Phi\) is an isomorphism.
	\end{proposition}
	
	\begin{proof}
		Nondegeneracy implies \(\Ker(\Phi)=\{0\}\), so \(\Phi\) is injective. Since \(\dim V=\dim V^*\), \(\Phi\) is an isomorphism.
	\end{proof}
	
	\begin{remark}
		This is the precise sense in which an inner product or nondegenerate bilinear form identifies \(V\) with \(V^*\).
	\end{remark}
	
	\section{Annihilators and Orthogonal Complements}
	
	\begin{definition}
		Let \(W\subseteq V\) be a subspace. Its \emph{annihilator} is
		\[
		W^\circ := \{f\in V^*: f(w)=0 \text{ for all } w\in W\}.
		\]
	\end{definition}
	
	This is the dual-space object naturally attached to a subspace \(W\).
	
	\begin{proposition}
		Under the isomorphism \(V\cong V^*\) induced by a nondegenerate bilinear form,
		\[
		W^\perp \quad\text{corresponds to}\quad W^\circ.
		\]
	\end{proposition}
	
	\begin{proof}
		A vector \(v\in V\) corresponds to the functional \(\varphi_v(x)=\ip{x}{v}\). This functional lies in \(W^\circ\) exactly when
		\[
		\varphi_v(w)=\ip{w}{v}=0 \quad \text{for all } w\in W,
		\]
		that is, exactly when \(v\in W^\perp\).
	\end{proof}
	
	Thus the orthogonal complement is the concrete realization inside \(V\) of the annihilator inside \(V^*\).
	
	\section{The Isomorphism \(V/W^\perp \cong W^*\)}
	
	Let \(W\subseteq V\) be a subspace, and assume \(V\) carries a nondegenerate bilinear form.
	
	Define
	\[
	\Psi:V\to W^*,\qquad \Psi(v)(w)=\ip{w}{v}.
	\]
	
	\begin{theorem}
		The map \(\Psi\) is linear and
		\[
		\Ker(\Psi)=W^\perp.
		\]
		Hence
		\[
		V/W^\perp \cong W^*.
		\]
	\end{theorem}
	
	\begin{proof}
		Linearity is immediate. Now
		\[
		v\in \Ker(\Psi)
		\iff \Psi(v)=0
		\iff \Psi(v)(w)=0 \text{ for all } w\in W
		\]
		\[
		\iff \ip{w}{v}=0 \text{ for all } w\in W
		\iff v\in W^\perp.
		\]
		By the First Isomorphism Theorem,
		\[
		V/\Ker(\Psi)\cong \im(\Psi).
		\]
		Since every functional on \(W\) arises from pairing with some vector in \(V\), the image is \(W^*\). Therefore
		\[
		V/W^\perp\cong W^*.
		\]
	\end{proof}
	
	\begin{remark}
		This theorem is one of the main bridges from linear algebra to coding theory.
	\end{remark}
	
	\section{Linear Codes and Dual Codes}
	
	\begin{definition}
		A \emph{linear code} of length \(n\) over \(\F_q\) is a subspace
		\[
		\C \subseteq \F_q^n.
		\]
		If \(\dim \C = k\), we say \(\C\) is an \([n,k]\)-code.
	\end{definition}
	
	\begin{definition}
		The \emph{dual code} of \(\C\) is
		\[
		\Cp := \{v\in \F_q^n : \ip{c}{v}=0 \text{ for all } c\in \C\}.
		\]
	\end{definition}
	
	\begin{remark}
		This is exactly the orthogonal complement of \(\C\) in \(\F_q^n\). The phrase ``dual code'' is coding-theoretic terminology for the orthogonal complement.
	\end{remark}
	
	\begin{theorem}
		For any linear code \(\C\subseteq \F_q^n\),
		\[
		\dim \C + \dim \Cp = n.
		\]
	\end{theorem}
	
	\begin{proof}
		This is the standard dimension formula for orthogonal complements under a nondegenerate bilinear form.
	\end{proof}
	
	\begin{remark}
		Over \(\mathbb{R}\) with a positive definite inner product, one always has
		\[
		V = W \oplus W^\perp.
		\]
		Over finite fields, the bilinear form is nondegenerate but need not be positive definite. Hence it is possible that
		\[
		\C \cap \Cp \neq \{0\}.
		\]
		So one must distinguish between \(\C+\Cp\) and \(\C\oplus\Cp\).
	\end{remark}
	
	\section{Why Coding Theory Calls \(\C^\perp\) the Dual Code}
	
	The term \emph{dual} is related to the passage from \(V\) to \(V^*\), but \(\Cp\) is not literally the same as \(\C^*\).
	
	\subsection{The distinction}
	
	\begin{itemize}[leftmargin=2em]
		\item \(\Cp\subseteq \F_q^n\) is a subspace of the ambient space.
		\item \(\C^*=\Hom(\C,\F_q)\) is the dual space of the code itself.
	\end{itemize}
	
	These are different objects.
	
	\subsection{The relation}
	
	The code \(\Cp\) corresponds to the annihilator \(\C^\circ\subseteq (\F_q^n)^*\). Because the standard bilinear form identifies \(\F_q^n\cong(\F_q^n)^*\), the annihilator appears concretely as a subspace of \(\F_q^n\), namely \(\Cp\).
	
	Thus the name ``dual code'' is justified because \(\Cp\) is the orthogonal-complement version of the annihilator of \(\C\).
	
	\section{Generator Matrices and Parity-Check Matrices}
	
	\begin{definition}
		A \emph{generator matrix} for an \([n,k]\)-code \(\C\) is a \(k\times n\) matrix \(G\) whose rows form a basis of \(\C\). Then
		\[
		\C = \{uG : u\in \F_q^k\}.
		\]
	\end{definition}
	
	\begin{definition}
		A \emph{parity-check matrix} for \(\C\) is a matrix \(H\) whose rows form a basis of \(\Cp\). Then
		\[
		c\in \C \iff Hc^T=0.
		\]
	\end{definition}
	
	\begin{proof}[Reason]
		Since the rows of \(H\) span \(\Cp\), the equation \(Hc^T=0\) exactly says that \(c\) is orthogonal to every vector in \(\Cp\), that is,
		\[
		c\in (\Cp)^\perp.
		\]
		Because the standard form on \(\F_q^n\) is nondegenerate,
		\[
		(\Cp)^\perp=\C.
		\]
	\end{proof}
	
	\begin{remark}
		So \(G\) describes the code by generators, while \(H\) describes the code by linear constraints.
	\end{remark}
	
	\section{The Dual Space \(\C^*\)}
	
	\begin{definition}
		The dual space of the code \(\C\) is
		\[
		\C^*:=\Hom(\C,\F_q).
		\]
	\end{definition}
	
	An element \(f\in \C^*\) is a linear functional on codewords:
	\[
	f:\C\to \F_q.
	\]
	
	If \(\dim \C=k\), then \(\dim \C^*=k\).
	
	\subsection{How a vector in \(\F_q^n\) produces an element of \(\C^*\)}
	
	Given \(v\in \F_q^n\), define
	\[
	f_v:\C\to \F_q,\qquad f_v(c)=\ip{c}{v}.
	\]
	This is a linear functional on \(\C\), so \(f_v\in \C^*\).
	
	Define
	\[
	\Phi:\F_q^n\to \C^*,\qquad \Phi(v)=f_v.
	\]
	
	\begin{theorem}
		\[
		\Ker(\Phi)=\Cp.
		\]
		Hence
		\[
		\F_q^n/\Cp \cong \C^*.
		\]
	\end{theorem}
	
	\begin{proof}
		By definition,
		\[
		v\in \Ker(\Phi)
		\iff \Phi(v)=0
		\iff f_v(c)=0 \text{ for all } c\in \C
		\]
		\[
		\iff \ip{c}{v}=0 \text{ for all } c\in \C
		\iff v\in \Cp.
		\]
		Now apply the First Isomorphism Theorem.
	\end{proof}
	
	\subsection{Interpretation}
	
	The quotient \(\F_q^n/\Cp\) identifies two ambient vectors \(v,v'\) exactly when they induce the same linear functional on \(\C\). Thus:
	
	\[
	v+\Cp
	\quad\leftrightarrow\quad
	\big(c\mapsto \ip{c}{v}\big).
	\]
	
	So \(\F_q^n/\Cp\) is the space of distinct linear observations of codewords coming from the ambient dot product.
	
	\section{The Dual Space \((\C^\perp)^*\)}
	
	Similarly, define
	\[
	(\Cp)^* := \Hom(\Cp,\F_q).
	\]
	Its elements are linear functionals on the dual code.
	
	Given \(y\in \F_q^n\), define
	\[
	g_y:\Cp\to \F_q,\qquad g_y(h)=\ip{h}{y}.
	\]
	This gives a linear map
	\[
	\Theta:\F_q^n\to (\Cp)^*,\qquad \Theta(y)=g_y.
	\]
	
	\begin{theorem}
		\[
		\Ker(\Theta)=\C.
		\]
		Hence
		\[
		\F_q^n/\C \cong (\Cp)^*.
		\]
	\end{theorem}
	
	\begin{proof}
		We have
		\[
		y\in \Ker(\Theta)
		\iff \ip{h}{y}=0 \text{ for all } h\in \Cp
		\iff y\in (\Cp)^\perp.
		\]
		Since the form is nondegenerate,
		\[
		(\Cp)^\perp=\C.
		\]
		Therefore \(\Ker(\Theta)=\C\), and the result follows by the First Isomorphism Theorem.
	\end{proof}
	
	\subsection{Interpretation}
	
	The quotient \(\F_q^n/\C\) identifies two ambient words \(y,y'\) exactly when they behave the same against all parity checks \(h\in\Cp\). So
	\[
	y+\C
	\quad\leftrightarrow\quad
	\big(h\mapsto \ip{h}{y}\big).
	\]
	
	This is the quotient that naturally appears in syndrome decoding.
	
	\section{Syndrome Decoding and the Quotient \(\F_q^n/\C\)}
	
	Let \(H\) be a parity-check matrix of \(\C\). For a received word \(y\in \F_q^n\), its \emph{syndrome} is
	\[
	s(y):=Hy^T.
	\]
	
	\begin{proposition}
		Two words \(y,y'\in \F_q^n\) have the same syndrome if and only if they lie in the same coset of \(\C\), that is,
		\[
		Hy^T=Hy'^T
		\iff y-y'\in \C
		\iff y+\C=y'+\C.
		\]
	\end{proposition}
	
	\begin{proof}
		\[
		Hy^T=Hy'^T
		\iff H(y-y')^T=0
		\iff y-y'\in \Ker(H)
		\iff y-y'\in \C.
		\]
	\end{proof}
	
	Thus syndromes classify cosets of \(\F_q^n/\C\).
	
	\subsection{Error-correction interpretation}
	
	Suppose \(c\in \C\) is sent and \(y=c+e\) is received. Then
	\[
	Hy^T = H(c+e)^T = Hc^T + He^T = 0 + He^T = He^T.
	\]
	So the syndrome depends only on the error vector \(e\), or equivalently only on the coset \(y+\C=e+\C\).
	
	This is why decoding is naturally phrased in terms of \(\F_q^n/\C\), not \(\F_q^n/\Cp\).
	
	\section{Summary of the Two Quotients}
	
	There are two natural quotient spaces associated to a code \(\C\subseteq \F_q^n\):
	
	\subsection*{First quotient}
	\[
	\F_q^n/\Cp \cong \C^*.
	\]
	Meaning: ambient vectors modulo those invisible to codewords.
	
	Interpretation: distinct linear functionals on the code.
	
	\subsection*{Second quotient}
	\[
	\F_q^n/\C \cong (\Cp)^*.
	\]
	Meaning: ambient words modulo adding a codeword.
	
	Interpretation: distinct syndrome classes; this is the quotient used in decoding.
	
	\subsection*{Memory aid}
	\[
	\boxed{
		\begin{aligned}
			\F_q^n/\Cp &\cong \C^*,\\
			\F_q^n/\C &\cong (\Cp)^*.
		\end{aligned}
	}
	\]
	
	\section{Direct Sum Issues: \(\C\oplus \Cp\)}
	
	Since
	\[
	\dim \C + \dim \Cp = n,
	\]
	one may ask whether
	\[
	\F_q^n = \C \oplus \Cp.
	\]
	
	This happens if and only if
	\[
	\C\cap \Cp = \{0\}.
	\]
	
	\begin{definition}
		A code \(\C\) is called \emph{LCD} (linear complementary dual) if
		\[
		\C\cap \Cp=\{0\}.
		\]
	\end{definition}
	
	Then and only then,
	\[
	\F_q^n=\C\oplus \Cp.
	\]
	
	\begin{remark}
		Over finite fields, it is possible that \(\C\cap \Cp\neq\{0\}\). For instance, a nonzero vector may be orthogonal to itself.
	\end{remark}
	
	\section{Worked Example over \(\F_2^3\)}
	
	Let
	\[
	\C = \Span\{(1,1,0),(0,1,1)\}\subseteq \F_2^3.
	\]
	Then \(\dim\C=2\), so \(\dim\Cp=1\).
	
	\subsection{A generator matrix}
	
	Take
	\[
	G=
	\begin{pmatrix}
		1&1&0\\
		0&1&1
	\end{pmatrix}.
	\]
	
	\subsection{Find the dual code}
	
	Let \(h=(a,b,c)\in \F_2^3\). Then \(h\in\Cp\) if and only if
	\[
	\ip{(1,1,0)}{(a,b,c)}=a+b=0,
	\]
	and
	\[
	\ip{(0,1,1)}{(a,b,c)}=b+c=0.
	\]
	Thus
	\[
	a=b=c.
	\]
	Hence
	\[
	\Cp=\{(0,0,0),(1,1,1)\}=\Span\{(1,1,1)\}.
	\]
	
	A parity-check matrix is therefore
	\[
	H=
	\begin{pmatrix}
		1&1&1
	\end{pmatrix}.
	\]
	
	\subsection{Interpretation of \(\C^*\)}
	
	Since \(\dim\C=2\), the space \(\C^*\) is 2-dimensional.
	
	For \(v\in \F_2^3\), the corresponding functional is
	\[
	f_v(c)=\ip{c}{v}.
	\]
	Two vectors \(v,v'\) define the same functional on \(\C\) iff \(v-v'\in\Cp\). Therefore
	\[
	\F_2^3/\Cp \cong \C^*.
	\]
	
	\subsection{Interpretation of \((\Cp)^*\)}
	
	Since \(\dim\Cp=1\), the space \((\Cp)^*\) is 1-dimensional.
	
	For \(y\in \F_2^3\), the corresponding functional on \(\Cp\) is
	\[
	g_y(h)=\ip{h}{y}.
	\]
	Two vectors \(y,y'\) define the same functional on \(\Cp\) iff \(y-y'\in\C\). Therefore
	\[
	\F_2^3/\C \cong (\Cp)^*.
	\]
	
	The syndrome of \(y\) is
	\[
	Hy^T = y_1+y_2+y_3.
	\]
	Thus there are exactly two cosets of \(\C\): one with syndrome \(0\), one with syndrome \(1\).
	
	\section{Conceptual Diagram}
	
	The entire story can be organized as follows:
	
	\[
	\begin{array}{ccc}
		\C & \subseteq & \F_q^n \\
		\cup && \cup \\
		\text{codewords} && \text{ambient words}
	\end{array}
	\]
	
	\[
	\Cp
	=
	\{h\in\F_q^n:\ip{c}{h}=0 \ \forall c\in\C\}
	\]
	
	\[
	\begin{aligned}
		G &: \text{rows span } \C,\\[4pt]
		H &: \text{rows span } \Cp.
	\end{aligned}
	\]
	
	\[
	\begin{aligned}
		\F_q^n/\Cp &\cong \C^*
		&& \text{(linear functionals on codewords)},\\[4pt]
		\F_q^n/\C &\cong (\Cp)^*
		&& \text{(linear functionals on parity checks / syndrome space)}.
	\end{aligned}
	\]
	
	\[
	c\in\C \iff Hc^T=0.
	\]
	
	\[
	y\mapsto Hy^T
	\quad\text{classifies cosets of }\F_q^n/\C.
	\]
	
	\section{Final Remarks}
	
	\begin{enumerate}[leftmargin=2em]
		\item The dual code \(\Cp\) is the orthogonal complement of \(\C\) inside \(\F_q^n\).
		\item The dual space \(\C^*\) is the space of linear functionals on \(\C\).
		\item These are different objects:
		\[
		\Cp \neq \C^*
		\]
		in general.
		\item However, they are linked by the bilinear form on \(\F_q^n\):
		\[
		\F_q^n/\Cp \cong \C^*.
		\]
		\item Dually,
		\[
		\F_q^n/\C \cong (\Cp)^*.
		\]
		\item In decoding, the quotient \(\F_q^n/\C\) is the natural one, because syndromes classify cosets modulo \(\C\).
		\item In functional or duality language, the quotient \(\F_q^n/\Cp\) is the natural one, because it records all ambient linear tests on codewords.
	\end{enumerate}

	\newpage
	\section{Algebraic Diagrams}
	
	We summarize the relations among the code, the dual code, the quotient spaces, and the dual spaces by commutative diagrams.
	
	\subsection{From the ambient space to \(\C^*\)}
	
	Define
	\[
	\Phi:\F_q^n\to \C^*,\qquad \Phi(v)(c)=\langle c,v\rangle.
	\]
	Its kernel is \(\C^\perp\). Hence:
	\[
	\F_q^n/\C^\perp \cong \C^*.
	\]
	
	\[
	\begin{tikzcd}[column sep=large]
		0 \arrow[r] & \C^\perp \arrow[r, hook] & \F_q^n \arrow[r, "\Phi"] & \C^* \arrow[r] & 0
	\end{tikzcd}
	\]
	
	Equivalently, the quotient map factors through the induced isomorphism:
	\[
	\begin{tikzcd}[column sep=huge]
		\F_q^n \arrow[r, two heads] \arrow[dr, "\Phi"'] & \F_q^n/\C^\perp \arrow[d, "\sim"] \\
		& \C^*
	\end{tikzcd}
	\]
	\subsection{From the ambient space to \((\C^\perp)^*\)}
	
	Define
	\[
	\Psi:\F_q^n\to (\C^\perp)^*,\qquad \Psi(y)(h)=\langle h,y\rangle.
	\]
	Its kernel is \(\C\). Hence:
	\[
	\F_q^n/\C \cong (\C^\perp)^*.
	\]
	
	\[
	\begin{tikzcd}[column sep=large]
		0 \arrow[r] & \C \arrow[r, hook] & \F_q^n \arrow[r, "\Psi"] & (\C^\perp)^* \arrow[r] & 0
	\end{tikzcd}
	\]
	
	Equivalently,
	\[
	\begin{tikzcd}[column sep=huge]
		\F_q^n \arrow[r, two heads] \arrow[dr, "\Psi"'] & \F_q^n/\C \arrow[d, "\sim"] \\
		& (\C^\perp)^*
	\end{tikzcd}
	\]
	
	\subsection{Generator and parity-check viewpoints}
	
	Let \(G\) be a generator matrix of \(\C\), and let \(H\) be a parity-check matrix of \(\C\). Then:
	
	\[
	\begin{tikzcd}[column sep=large, row sep=large]
		\F_q^k \arrow[r, "u \mapsto uG"] & \C \arrow[r, hook] & \F_q^n \arrow[r, "y \mapsto Hy^T"] & \F_q^{n-k}
	\end{tikzcd}
	\]
	
	Here:
	
	\begin{itemize}
		\item the map \(u \mapsto uG\) parametrizes codewords;
		\item the inclusion \(\C \hookrightarrow \F_q^n\) places the code inside the ambient space;
		\item the map \(y \mapsto Hy^T\) gives the syndrome;
		\item \(\ker(H)=\C\).
	\end{itemize}

	\newpage
	\subsection{Duality square}
	
	\[
	\begin{tikzcd}[row sep=large, column sep=huge]
		\F_q^n/\C^\perp \arrow[r, "\sim"] & \C^* \\
		\F_q^n/\C \arrow[r, "\sim"] \arrow[u, dashed] & (\C^\perp)^* \arrow[u, dashed]
	\end{tikzcd}
	\]
	
	\section{Algebraic Diagrams via \texttt{tikz-cd}}
	
	The two fundamental quotient identifications in coding theory are
	\[
	\F_q^n/\C^\perp \cong \C^*,
	\qquad
	\F_q^n/\C \cong (\C^\perp)^*.
	\]
	
	The first comes from the pairing of ambient vectors with codewords:
	\[
	\Phi:\F_q^n\to \C^*,\qquad \Phi(v)(c)=\langle c,v\rangle.
	\]
	Its kernel is \(\C^\perp\), so we obtain the short exact sequence
	\[
	\begin{tikzcd}[column sep=large]
		0 \arrow[r] & \C^\perp \arrow[r, hook] & \F_q^n \arrow[r, "\Phi"] & \C^* \arrow[r] & 0.
	\end{tikzcd}
	\]
	Therefore,
	\[
	\begin{tikzcd}[column sep=huge]
		\F_q^n \arrow[r, two heads] \arrow[dr, "\Phi"'] & \F_q^n/\C^\perp \arrow[d, "\sim"] \\
		& \C^*.
	\end{tikzcd}
	\]
	
	The second comes from the pairing of ambient vectors with parity checks:
	\[
	\Psi:\F_q^n\to (\C^\perp)^*,\qquad \Psi(y)(h)=\langle h,y\rangle.
	\]
	Its kernel is \(\C\), so we obtain
	\[
	\begin{tikzcd}[column sep=large]
		0 \arrow[r] & \C \arrow[r, hook] & \F_q^n \arrow[r, "\Psi"] & (\C^\perp)^* \arrow[r] & 0.
	\end{tikzcd}
	\]
	Therefore,
	\[
	\begin{tikzcd}[column sep=huge]
		\F_q^n \arrow[r, two heads] \arrow[dr, "\Psi"'] & \F_q^n/\C \arrow[d, "\sim"] \\
		& (\C^\perp)^*.
	\end{tikzcd}
	\]
	
	These diagrams express the two complementary viewpoints:
	
	\begin{itemize}
		\item \(\F_q^n/\C^\perp \cong \C^*\): ambient vectors modulo those invisible to codewords;
		\item \(\F_q^n/\C \cong (\C^\perp)^*\): ambient vectors modulo codewords, which is the quotient relevant for syndrome decoding.
	\end{itemize}

	\newpage
%	\begin{table}[h]
%		\centering
%		\renewcommand{\arraystretch}{1.35}
%		\begin{tabular}{|p{3cm}|p{2cm}|p{4.6cm}|p{3.8cm}|}
%			\hline
%			\textbf{Concept} & \textbf{Meaning} & \textbf{Formal condition} & \textbf{Coding-theory consequence} \\
%			\hline
%			Nondegenerate bilinear form
%			&
%			No nonzero vector is orthogonal to \emph{every} vector in the whole ambient space.
%			&
%			For \(v\in \F_q^n\),
%			$
%			\langle x,v\rangle=0 \ \forall x\in \F_q^n
%			\quad\Longrightarrow\quad
%			v=0.
%			$
%			&
%			Ensures the orthogonal-complement operation behaves well, e.g.
%			$
%			\dim \mathcal C+\dim \mathcal C^\perp=n,
%			\qquad
%			(\mathcal C^\perp)^\perp=\mathcal C.
%			$
%			\\
%			\hline
%			Positive-definite form
%			&
%			A vector can be orthogonal to itself only if it is zero.
%			&
%			$
%			\langle v,v\rangle=0 \Longrightarrow v=0.
%			$
%			&
%			This usually fails over finite fields. Hence one cannot generally conclude
%			$
%			\mathcal C\cap \mathcal C^\perp=\{0\}.
%			$
%			\\
%			\hline
%			Isotropic vector
%			&
%			A nonzero vector orthogonal to itself.
%			&
%			$
%			v\neq 0,
%			\qquad
%			\langle v,v\rangle=0.
%			$
%			&
%			Makes self-orthogonality possible. Such vectors can lie in both
%			\(\mathcal C\) and \(\mathcal C^\perp\).
%			\\
%			\hline
%			Self-orthogonal code
%			&
%			Every codeword is orthogonal to every codeword of the code.
%			&
%			$
%			\mathcal C\subseteq \mathcal C^\perp.
%			$
%			&
%			All codewords satisfy all parity checks coming from the code itself.
%			In particular,
%			$
%			\mathcal C\cap \mathcal C^\perp=\mathcal C.
%			$
%			\\
%			\hline
%			Self-dual code
%			&
%			The code equals its dual.
%			&
%			$
%			\mathcal C=\mathcal C^\perp.
%			$
%			&
%			The code is exactly its own parity-check space. Necessarily
%			$
%			\dim \mathcal C=\frac n2.
%			$
%			\\
%			\hline
%			LCD code
%			&
%			The code and its dual intersect trivially.
%			&
%			\[
%			\mathcal C\cap \mathcal C^\perp=\{0\}.
%			\]
%			&
%			Then the ambient space splits as a direct sum:
%			\[
%			\F_q^n=\mathcal C\oplus \mathcal C^\perp.
%			\]
%			\\
%			\hline
%		\end{tabular}
%		\caption{Comparison of nondegeneracy, positive-definiteness, isotropic vectors, self-orthogonal codes, self-dual codes, and LCD codes.}
%	\end{table}
	
	\subsection*{Nondegenerate bilinear form}
	
	A bilinear form \(\langle\ ,\ \rangle\) on \(\F_q^n\) is called \emph{nondegenerate} if
	\[
	\langle x,v\rangle=0 \quad \text{for all } x\in \F_q^n
	\quad \Longrightarrow \quad
	v=0.
	\]
	This means that no nonzero vector is orthogonal to every vector in the whole ambient space.
	
	In coding theory, this property ensures that orthogonal complements behave well. In particular, for any linear code \(\mathcal C\subseteq \F_q^n\),
	\[
	\dim \mathcal C+\dim \mathcal C^\perp=n,
	\qquad
	(\mathcal C^\perp)^\perp=\mathcal C.
	\]
	
	\subsection*{Positive-definite form}
	
	A bilinear form is called \emph{positive-definite} in the Euclidean sense if
	\[
	\langle v,v\rangle=0 \quad \Longrightarrow \quad v=0.
	\]
	This is the property that makes orthogonality in \(\mathbb R^n\) so rigid.
	
	However, over finite fields the standard bilinear form is generally \emph{not} positive-definite. Therefore it may happen that a nonzero vector satisfies
	\[
	\langle v,v\rangle=0.
	\]
	This is the main reason why one cannot generally conclude that
	\[
	\mathcal C\cap \mathcal C^\perp=\{0\}.
	\]
	
	\subsection*{Isotropic vector}
	
	A nonzero vector \(v\in \F_q^n\) is called \emph{isotropic} if
	\[
	\langle v,v\rangle=0.
	\]
	So an isotropic vector is a nonzero vector that is orthogonal to itself.
	
	Such vectors are impossible in a positive-definite space, but they often exist over finite fields. Their existence is exactly what allows a codeword to lie in both \(\mathcal C\) and \(\mathcal C^\perp\).
	
	\subsection*{Self-orthogonal code}
	
	A linear code \(\mathcal C\subseteq \F_q^n\) is called \emph{self-orthogonal} if
	\[
	\mathcal C\subseteq \mathcal C^\perp.
	\]
	This means that every codeword is orthogonal to every codeword of the code, including itself.
	
	In that case, every vector in \(\mathcal C\) automatically satisfies all parity checks coming from \(\mathcal C\) itself. In particular,
	\[
	\mathcal C\cap \mathcal C^\perp=\mathcal C.
	\]
	
	\subsection*{Self-dual code}
	
	A linear code \(\mathcal C\subseteq \F_q^n\) is called \emph{self-dual} if
	\[
	\mathcal C=\mathcal C^\perp.
	\]
	So the code coincides exactly with its dual code.
	
	This means that the code is simultaneously the space of codewords and the space of all parity checks for those codewords. Since
	\[
	\dim \mathcal C+\dim \mathcal C^\perp=n,
	\]
	a self-dual code must satisfy
	\[
	\dim \mathcal C=\frac n2.
	\]
	
	\subsection*{LCD code}
	
	A linear code \(\mathcal C\subseteq \F_q^n\) is called \emph{LCD} (linear complementary dual) if
	\[
	\mathcal C\cap \mathcal C^\perp=\{0\}.
	\]
	This means that the code and its dual meet only at the zero vector.
	
	In that case the ambient space splits as a direct sum:
	\[
	\F_q^n=\mathcal C\oplus \mathcal C^\perp.
	\]
	So every vector in \(\F_q^n\) can be written uniquely as the sum of one vector from \(\mathcal C\) and one vector from \(\mathcal C^\perp\).
	
	\subsection*{Main point}
	
	The key distinction is the following:
	
	\begin{itemize}
		\item nondegeneracy says that no nonzero vector is orthogonal to the entire ambient space;
		\item positive-definiteness says that no nonzero vector is orthogonal to itself.
	\end{itemize}
	
	Over finite fields, the standard bilinear form is nondegenerate but usually not positive-definite. Hence isotropic vectors may exist, and therefore it is possible that
	\[
	\mathcal C\cap \mathcal C^\perp\neq \{0\}.
	\]
	
	\subsection*{Example}
	
	Over \(\F_2\), let
	\[
	\mathcal C=\operatorname{span}\{(1,1)\}\subseteq \F_2^2.
	\]
	Then
	\[
	\langle (1,1),(1,1)\rangle = 1+1 = 0 \pmod 2.
	\]
	So \((1,1)\) is isotropic. Hence
	\[
	\mathcal C=\mathcal C^\perp,
	\]
	and therefore
	\[
	\mathcal C\cap \mathcal C^\perp=\mathcal C\neq \{0\}.
	\]
	This shows that trivial intersection does not hold in general over finite fields.


	\begin{example}
		Over \(\F_2\), let
		\[
		\mathcal C=\operatorname{span}\{(1,1)\}\subseteq \F_2^2.
		\]
		Then
		\[
		\langle (1,1),(1,1)\rangle = 1+1 = 0 \pmod 2.
		\]
		So \((1,1)\) is isotropic, and hence
		\[
		\mathcal C=\mathcal C^\perp.
		\]
		Therefore
		\[
		\mathcal C\cap \mathcal C^\perp=\mathcal C\neq \{0\}.
		\]
		This shows that trivial intersection does not hold in general over finite fields.
	\end{example}
\end{document}